\documentclass[../main.tex]{subfiles}
\begin{document}

% Target length: >= 5 pages.

Chapter~\ref{chapter:Experiment} presented the system as a whole.
This chapter isolates its six contributions and examines each as a problem, a solution, and a result.
Table~\ref{table:problems} lists the problem each one answers.
Five of them concern building a reliable and maintainable \gls{LLM} application around an external model; the last returns to the domain problem of contradictory official data.

\begin{table}[H]
    \centering
    \caption{Problems addressed by the system's contributions}
    \label{table:problems}
    {\footnotesize
    \setlength{\tabcolsep}{5pt}
    \begin{tabular}{|p{4.5cm}|p{6.3cm}|c|}
        \hline
        \textbf{Problem} & \textbf{Why it matters} & \textbf{Addressed in} \\
        \hline
        The hosted \gls{LLM} fails as a hard outage or as an unusable (malformed) response & A naive caller crashes a consultation instead of degrading & Section~\ref{section:5.1} \\
        \hline
        Sources share no format and some \gls{PDF}s are scanned images & A uniform pipeline cannot ingest them, leaving the store empty & Section~\ref{section:5.2} \\
        \hline
        A multi-stage run is opaque & There is no single place to see why a run is slow, costly, or wrong & Section~\ref{section:5.3} \\
        \hline
        Orchestration outside the advisory graph was imperative and the knowledge flow was duplicated & Control flow drifts and cannot be traced uniformly & Section~\ref{section:5.4} \\
        \hline
        Repeated factual questions re-run the full retrieval and generation & Inference quota and latency are spent re-answering settled questions, and a naive cache serves stale answers & Section~\ref{section:5.5} \\
        \hline
        Official sources contradict each other on quotas and cutoff scores & A tool that silently picks one value is often wrong without saying so & Section~\ref{section:5.6} \\
        \hline
    \end{tabular}}
\end{table}

\section{Resilient LLM inference gateway}
\label{section:5.1}

\textbf{Problem.}
The assistant depends on a hosted \gls{LLM} reached over the network, and that dependency fails in two ways.
A \textit{hard failure} means the \gls{API} is unreachable, unauthenticated, or rate-limited, which is common because the inference quota is finite and shared across every reasoning step.
A \textit{soft failure} means the call succeeds but the body cannot be used: the model wraps its \gls{JSON} in prose, emits broken syntax, or returns a value outside the expected set.
A naive caller that trusts every \gls{LLM} call would let either failure crash a consultation, breaking the graceful-degradation requirement NFR-01.
The two cases also need opposite responses: retrying a rate-limited endpoint is pointless, while retrying a malformed response can succeed.

\textbf{Solution.}
All model access goes through one gateway (Section~\ref{section:4.1}) that makes the two failure modes first-class; its control flow is given in Algorithm~\ref{alg:gateway}.
A hard failure is raised as an \texttt{InferenceError}; a soft failure is reported on the result as a \texttt{failure\_type} of \texttt{STRUCTURE\_FAILURE} rather than as an exception.
A model registry gives each agent a policy: a primary model, an optional fallback model, a retry budget, and an \texttt{allow\_fallback} flag that defaults to off.
On a hard error the gateway stops retrying the same model at once and moves to the fallback, because retrying an unreachable endpoint wastes time; on a structure failure it retries within the budget and then falls back.
With no fallback configured, it re-raises the hard error so the call site can degrade, and every call site does exactly that, wrapping the call and substituting deterministic output.
This lets a cheap, high-volume step like intent classification run with no fallback and drop to deterministic keywords, while a step whose output is harder to rebuild gets a backup model.
Every attempt is logged with its model, attempt number, \texttt{failure\_type}, and a \texttt{used\_fallback} flag, so a degraded call is never silent.
The clearest case is intent routing: when the model is unavailable, a keyword classifier mirrors the prompt's own priority, checking knowledge-topic keywords before advisory phrases, so a factual question is never misrouted into a full advisory run, and unrecognized messages ask for clarification rather than silently re-run (commit \texttt{c2ef582}).

\begin{algorithm}[H]
\caption{Inference gateway: hard- vs structure-failure handling}
\label{alg:gateway}
\KwIn{request $q$ with an agent name}
\KwOut{a usable \texttt{InferenceResult}, or a raised \texttt{InferenceError} for the call site to handle}
$p \leftarrow$ registry policy for $q$'s agent \tcp*[r]{primary, fallback, retries, allow\_fallback}
$e \leftarrow \varnothing$ \tcp*[r]{remembered hard error}
\For{$\text{attempt} \leftarrow 0$ \KwTo $p.\textit{maxRetries}$}{
  $r \leftarrow \textsc{Generate}(q, p.\textit{primaryModel})$\;
  \If{\textsc{Generate} raised \texttt{InferenceError} $\varepsilon$ \tcp*[r]{hard: network/auth/rate-limit}}{
    $e \leftarrow \varepsilon$; record \textsc{ApiError}; \textbf{break} \tcp*[r]{stop retrying this model}
  }
  record $r.\textit{failureType}$\;
  \lIf{$r.\textit{failureType} \neq \textsc{StructureFailure}$}{\Return $r$}
}
\eIf{$p.\textit{allowFallback}$ \textbf{and} $p.\textit{fallbackModel}$ exists}{
  $r \leftarrow \textsc{Generate}(q, p.\textit{fallbackModel})$ \tcp*[r]{re-raises on hard error}
  \Return $r$\;
}{
  \lIf{$e \neq \varnothing$}{\textbf{raise} $e$ \tcp*[r]{let the call site degrade}}
  \Return $r$ \tcp*[r]{last structure-failure result}
}
\end{algorithm}

\textbf{Result.}
The assistant stays usable through model outages.
Intent routing keeps working with no model call, knowledge answering degrades to an honest ``no data'' reply rather than a fabricated one, and malformed outputs are absorbed by retry and fallback instead of surfacing as crashes.
One gateway means every reasoning step - routing, conflict tie-breaking, knowledge answering, and \gls{OCR} - inherits the same resilience, and the telemetry shows when a fallback or a structure failure occurred, so an outage is visible rather than silent.

\section{End-to-end ingestion of heterogeneous official sources}
\label{section:5.2}

\textbf{Problem.}
The data that makes the assistant useful lives in sources with no common format: each university structures its admission pages differently, proposal documents are published as \gls{PDF}s, and some of those \gls{PDF}s are scanned images with no text layer.
A pipeline that assumes clean, machine-readable input cannot ingest this material, yet without it the canonical store and the knowledge corpus stay empty.
Scanned documents are the hardest case, because recovering their text needs \gls{OCR}, and model-based \gls{OCR} brings a failure mode of its own.

\textbf{Solution.}
Ingestion is a staged, configurable pipeline (Section~\ref{section:4.1}): per-source configuration picks the fetch strategy, a router predicts each document's type, and a matching parser and extractor turn it into raw facts that are then normalized against per-school dictionaries into canonical records.
Because the configuration and dictionaries are data rather than code, adding a school is mostly a matter of configuration.
Scanned \gls{PDF}s go through a hybrid per-page extractor (Section~\ref{section:3.4}, Algorithm~\ref{alg:ocr}): a page whose text layer is at least 50 characters long is read directly, and only shorter, effectively image-only pages are rendered to PNG at roughly 200 \gls{DPI} and sent to the multimodal gateway, keeping cost proportional to the number of scanned pages.
The key part is the safeguard against \gls{OCR} repetition loops (commit \texttt{b41dd9f}).
On a page with merged table cells the model can lock into emitting one character endlessly - one NEU page produced on the order of a million dash characters.
The extractor flags an output as \textit{degenerate} when it passes 20{,}000 characters, or when one character makes up more than 80\% of at least 2{,}000 non-whitespace characters, and retries the page once at temperature 0.3, up from 0.0, to break the loop.
If it is still degenerate, the page is marked failed and the rest of the document is ingested normally, so one bad page cannot discard a whole document.
A document that yields no text at all is left un-ingested rather than recorded as empty.

\begin{algorithm}[H]
\caption{Hybrid per-page extraction with degenerate-OCR detection}
\label{alg:ocr}
\KwIn{a scanned or mixed \gls{PDF}}
\KwOut{per-page text, each tagged \texttt{text\_layer}, \texttt{ocr}, or \texttt{failed}}
\ForEach{page in document}{
  $b \leftarrow$ stripped pdfplumber text layer of the page\;
  \lIf{$|b| \geq 50$}{record $(b, \texttt{text\_layer})$; \textbf{continue}}
  $\textit{img} \leftarrow$ render the page to PNG at $\approx 200$ \gls{DPI}\;
  \For{$\tau \in (0.0,\ 0.3)$ \tcp*[r]{retry once at raised temperature}}{
    $t \leftarrow \textsc{GatewayOCR}(\textit{img},\ \tau)$\;
    \lIf{\textbf{not} \textsc{Degenerate}$(t)$}{record $(t, \texttt{ocr})$; \textbf{continue to next page}}
  }
  record $(\texttt{""}, \texttt{failed})$ \tcp*[r]{both attempts degenerate}
}
\lIf{every page is empty}{\textbf{raise} - leave the document un-ingested}
\BlankLine
\SetKwProg{Fn}{Function}{:}{}
\Fn{\textsc{Degenerate}$(t)$}{
  \lIf{$|t| > 20000$}{\Return true}
  $c \leftarrow t$ with all whitespace removed\;
  \lIf{$|c| \geq 2000$ \textbf{and} $\max_x \text{count}(x, c) / |c| > 0.8$}{\Return true}
  \Return false\;
}
\end{algorithm}

Two parts of this control flow are deliberate.
First, the jump from the cheap text layer to the expensive vision call happens only when needed: the renderer is loaded lazily and used only for pages below the 50-character threshold, so a born-digital document never loads the rasterizer, and the number of paid \gls{OCR} calls stays proportional to the count of scanned pages.
Second, the document-level guard separates a partly failed document from an empty one for a concrete reason: ingested documents are deduplicated by content hash, so recording a zero-text document as ingested would remember its hash and skip it forever, hiding the failure; raising instead leaves the document outside the corpus, where a later run can retry it.

\textbf{Result.}
The pipeline ingests both structured and unstructured official sources end to end.
The canonical store holds 136 admission records and 699 cutoff records for HUST and 20 admission records and 16 cutoff records for VNU-UET, and the knowledge corpus holds 93 documents and 692 chunks drawn from HUST, NEU, VNU-UET, and national \gls{MOET} regulations, as reported in Section~\ref{section:4.3}.
The degenerate-output safeguard is what makes the scanned part of that corpus ingestible: without it, one repetition-loop page would either bloat the store with junk or abort the document; with it, the bad page is isolated and the rest is kept.

\section{Unified observability over advisory runs}
\label{section:5.3}

\textbf{Problem.}
An advisory run is opaque from the outside: one consultation fans out into six pipeline stages and several model calls, some of which retry or fall back, and when a run is slow, expensive, or wrong there is no single place to see why.
An earlier version wrote per-stage events to its own database table and showed them in an in-app panel (Section~\ref{section:3.6}), but that view showed stage timings only - not the prompt, the raw response, the token usage, or the per-attempt retry and fallback behavior, which is exactly what is needed to diagnose a degraded or costly run.
Maintaining a custom trace store and viewer also cost code and tests for a capability that purpose-built tools already provide.

\textbf{Solution.}
I retired the custom trace store and its in-app panel in favor of a single observability sink, Langfuse, integrated through a dedicated \texttt{observability} package.
A run opens one root span keyed by a trace identifier derived from the run's \texttt{trace\_run\_id}, so the external trace and the database run cross-reference by the same integer, and the span's session identifier is the chat session token, so every run of one conversation groups together.
The same \texttt{traced} decorator that wraps each advisory node (Section~\ref{section:4.1}) opens a child span per stage, so the fixed graph of Section~\ref{section:3.3} yields a trace tree of the same shape on every run.
Each model call is recorded as a generation observation carrying the model name, the raw prompts, the raw output, the token usage, the latency, the attempt number, the \texttt{used\_fallback} flag, and the \texttt{failure\_type} of Section~\ref{section:5.1}.
Because each retry and fallback is a separate generation, the gateway's resilience becomes directly visible.
Capturing token usage needed one small change: the Gemini provider, which previously kept only the response text, now reads the response's usage metadata into an optional \texttt{usage} field.
The integration keeps the graceful-degradation discipline of the rest of the system: the client is disabled by default and built lazily, a missing client makes every helper a no-op, and any backend error is caught and logged, so observability can never become a new failure mode.

\textbf{Result.}
Every advisory run is now a single queryable trace - a tree of six stage spans with the model calls recorded beneath the stages that make them - that exposes the prompts, outputs, token usage, latency, and fallback path the earlier panel could not.
One external sink replaced the parallel trace store and viewer while adding the per-call and cost visibility that was missing, and self-hosting keeps the captured prompts and student messages on the same machine as the rest of the system, consistent with the anonymity posture of NFR-03.

\section{Declarative orchestration of the conversation and knowledge flows}
\label{section:5.4}

\textbf{Problem.}
The advisory pipeline was a declarative graph from the start (Section~\ref{section:3.3}), but the code around it was not.
A conversation turn was a hard-coded sequence: extract the profile, run deterministic guards for reset, continuation, and correction, classify intent, then branch into one of several handlers, with the control flow buried in nested conditionals.
The knowledge flow had a second problem: the same embed-retrieve-augment-gate-generate sequence was reached from three call sites (the inline turn, the batch fan-out, and the hybrid comparison), inviting divergence and giving those paths no shared place to be traced.

\textbf{Solution.}
I expressed the conversation turn and the knowledge flow as LangGraph graphs, without adding any run-time model autonomy, so the determinism of Section~\ref{section:3.3} is preserved.
The turn graph makes the guards conditional-edge nodes and the intent router a routing node that sends each message to a handler - enqueue an advisory run, answer a knowledge question, run a hybrid comparison, hold a conversational reply, ask for clarification, or decline an out-of-scope request - so the routes are now graph edges rather than hidden branches.
The knowledge flow is defined once as a reusable subgraph (embed, retrieve school-scoped chunks, augment with national chunks, gate on confidence, generate) behind the knowledge service's unchanged \texttt{answer} method, so all three callers share one definition and the confidence gate of Algorithm~\ref{alg:retrieval} cannot drift between them.
The advisory graph itself is unchanged and still runs on the background worker; the turn graph only enqueues advisory runs, keeping the durable-queue design of Section~\ref{section:4.2} intact.
Because every node uses the same tracing decorator, the unified observability of Section~\ref{section:5.3} extends across all three flows.
The refactoring was guarded by characterization tests that pin the turn outcome for every route before the structure changed, so the behavior is preserved rather than just re-expressed.

\textbf{Result.}
The control flow is now declarative end to end: the per-turn routing, the knowledge subgraph, and the advisory pipeline are all graphs over typed state, and the knowledge subgraph is defined once and reused by the inline, fan-out, and hybrid paths.
The same node wrapping extends the observability of Section~\ref{section:5.3} to the synchronous turn and knowledge paths the run-only tracer missed, and the full test suite of Section~\ref{section:4.4} passing after the change is the evidence that the externally observable behavior is unchanged.

\section{Version-gated semantic answer cache}
\label{section:5.5}

\textbf{Problem.}
Factual questions cluster heavily: many students ask the same few things (``how much is tuition?'', ``what majors does this school offer?''), and answering each one re-embeds the question, runs a vector search, and calls the model.
Repeating that work for an already-answered question spends inference quota and adds latency for no new information.
A naive cache is unsafe here, because the knowledge corpus is re-ingested during the admission season (NFR-02), so a cached answer can outlive the document it came from and become wrong.

\textbf{Solution.}
The knowledge service is fronted by a semantic cache (migration \texttt{019}, Section~\ref{section:4.2}) that reuses a past answer when a new question is close enough, while gating every reuse on freshness.
On a miss, the answer is stored together with the question's 768-dimensional embedding, its confidence, and a stamp recording the version of its data scope, indexed for vector search with the same \gls{HNSW} machinery as the main corpus (Section~\ref{section:3.2}).
On a new question, the cache embeds it, finds the nearest stored entry within the same school-and-topic scope, and accepts it only when the similarity clears a threshold \textit{and} the stored scope version still equals the current one; otherwise it is a miss and the question flows through to the subgraph of Section~\ref{section:5.4}.
Every ingest bumps a scope's version, which invalidates exactly the cached answers whose data may have changed.
Like every other external-dependency feature, the cache is disabled by default; when it is off the service behaves exactly as before, and the result carries a flag indicating whether it came from the cache.

\textbf{Result.}
When enabled, a repeated factual question is answered from the cache without a model call, and a re-ingest of the relevant scope forces the next identical question to be recomputed, so the speed-up never comes at the cost of a stale answer.
The hit/miss path, the similarity gate, and stale-invalidation-on-ingest are covered by integration tests in the suite of Section~\ref{section:4.4}; a measured hit rate is left to a production deployment, since it depends on a live question stream the development corpus does not have.

\section{Conflict-aware data consolidation}
\label{section:5.6}

\textbf{Problem.}
Official admission data is not only fragmented but inconsistent: two authoritative publications can state different quotas for the same program, and two sources can record different historical cutoff scores for the same program and year.
As Chapter~\ref{chapter:Related_works} showed, every existing channel resolves such disagreement silently - it picks one value or shows none - which hides exactly the uncertainty a student needs before applying.
Consolidating these sources therefore cannot mean overwriting one value with another; the disagreement must survive into the recommendation.
Such contradictions are rare in the corpus gathered so far, so this is best seen as a safeguard for rare but high-impact cases: its value is what it prevents when it fires, not how often.
The data-fusion literature frames the options as resolving, ignoring, or deferring a conflict~\cite{bleiholder2009fusion}, and truth-discovery work shows that no source can be trusted unconditionally~\cite{li2016truth}; both motivate a conservative stance, resolving only on an objective, auditable signal and otherwise deferring.

\textbf{Solution.}
Conflict awareness is built in three layers, summarized in Algorithm~\ref{alg:conflict}.
At the storage layer (Section~\ref{section:4.2}), the canonical tables are unique per source, so contradictory publications coexist as separate rows rather than overwriting one another.
At the detection layer, the conflict service groups records by program, year, and method - the conflict key - and flags a conflict whenever a group holds two or more distinct values; quota and cutoff conflicts are detected separately, because they have different keys and rules.
At the resolution layer, the two kinds are treated differently.
A quota conflict is settled by a deterministic ranking on a fixed priority order: trust level first, then agreement across sources, then recency, then source confidence.
If the top option outranks its nearest differing challenger on any one of these, it is accepted; if every criterion ties, the field is marked uncertain and every competing value is shown.
No language model is consulted - the conflict agent wires only these comparators - so the resolution is fully auditable, and a fixed order rather than a single weighted score keeps the deciding criterion recoverable.
The agreement criterion counts how many independent sources report a value, the practical form of the truth-discovery idea that agreement is itself evidence~\cite{li2016truth}: a value asserted by two publications outranks an equally trusted value asserted by one.
A cutoff conflict is not resolved by ranking.
Instead the system checks whether the disagreement is \textit{decision-changing} for this student - whether the competing cutoff values would place the student's score into different fit categories.
If it is, the field is left uncertain and every competing value is shown; if it is not, or if the student's score or method is not on a comparable scale, the system references the most trusted value while still listing all of them.

\begin{algorithm}[H]
\caption{Conflict detection and resolution (per recommended program)}
\label{alg:conflict}
\KwIn{candidate records $R$ for a program; student score $s$ and method $m$}
\KwOut{a resolved value, or an uncertainty annotation listing all competing values}
Group $R$ by the conflict key $(\textit{program}, \textit{year}, \textit{method})$, separately for quota and cutoff fields\;
\ForEach{group $g$ with $\geq 2$ distinct quota values}{
  rank the options of $g$ by $(\textit{trust}, \textit{agreement}, \textit{recency}, \textit{confidence})$, descending\;
  \eIf{top option strictly outranks its nearest differing challenger on some axis}{
    resolve quota to the top-ranked value\;
  }{
    mark quota \textbf{uncertain}; surface every competing value\;
  }
}
\ForEach{group $g$ with $\geq 2$ distinct cutoff scores}{
  order the options of $g$ by trust level (then score)\;
  \eIf{$s$ known \textbf{and} $m$ on a comparable scale}{
    $F \leftarrow \{\,\textsc{ClassifyFit}(s, v) : v \in \textit{values}(g)\,\}$\;
    \eIf{$|F| > 1$ \tcp*[r]{decision-changing}}{
      mark cutoff \textbf{uncertain}; list every competing value\;
    }{
      reference the most-trusted value; still list all values\;
    }
  }{
    reference the most-trusted value; still list all values\;
  }
}
\end{algorithm}

\textbf{Result.}
The pipeline surfaces conflicts instead of hiding them.
When a program's sources disagree in a way that matters, the recommendation marks the field uncertain, lists the differing values, and advises the student to verify against the latest official announcement, the behavior required by edge case EC-16 and the conflict-disclosure use case of Chapter~\ref{chapter:Related_works}.
EC-16 and EC-17 are among the 17 satisfied cases reported in Section~\ref{section:4.4}.
Because real contradictions are rare in the live corpus, the subsystem is also tested with controlled synthetic conflicts (Section~\ref{subsection:4.4.6}): across five injected disagreements and three agreeing controls it detects every conflict with no false positives, resolves the cases that have an objective winner, and - the result that matters most - leaves an equal-trust tie and a decision-changing cutoff disagreement unresolved and surfaced rather than silently settled.

This chapter examined the system's six contributions.
Five make the application reliable and maintainable: a resilient inference gateway that separates hard from soft model failures, end-to-end ingestion of heterogeneous and scanned official sources with a safeguard against \gls{OCR} failure, unified observability that turns every run into one queryable trace, declarative orchestration of the conversation and knowledge flows as reusable graphs, and a version-gated cache that reuses answers without serving stale ones.
The last, conflict-aware consolidation, preserves and discloses contradictory source data, a safeguard for the rare but high-impact case where official sources disagree.
Chapter~\ref{chapter:conclusion} concludes the thesis and sets out the remaining work.

\end{document}
